\begingroup
\renewcommand{\arraystretch}{1.1}
\begin{table}[H]
	\centering
	\begin{tabular}{>{\bfseries}l >{\raggedright\arraybackslash}p{0.72\textwidth}}
		\hline\hline
		BAYGAUD & BAYesian GAUssian Decomposer\\
		BEARS & 25-Beam Array Receiver System\\
		CNM & Milieu neutre froid (\textit{cold neutral medium})\\
		CO & Monoxyde de carbone\\
		FUV & Ultraviolet lointain (\textit{far-ultraviolet})\\
		GALEX NGS & Galaxy Evolution Explorer Nearby Galaxies Survey\\
		GMCs & Nuages moléculaires géants (\textit{giant molecular clouds})\\
		H$_2$ & Hydrogène moléculaire\\
		HERACLES & HERA CO-Line Extragalactic Survey\\
		HI & Hydrogène atomique neutre\\
		HII & Hydrogène ionisé, région H II\\
		HIM & Milieu ionisé chaud (\textit{hot ionized medium})\\
		IR & Infrarouge \\
		IRAM & Institut de Radioastronomie Millimétrique\\
		MIS & Milieu interstellaire \\
		ISRF & Champ de rayonnement interstellaire (\textit{interstellar radiation field})\\
		LNM & Milieu neutre de température intermédiaire (\textit{lukewarm neutral medium})\\
		NRAO & National Radio Astronomy Observatory\\
		%NRO & Nobeyama Radio Observatory\\
		ROHSA & Regularized Optimization for Hyper-Spectral Analysis\\
		S$^4$G & Spitzer Survey of Stellar Structure in Galaxies\\
		SFE & Efficacité de formation stellaire (\textit{star formation efficiency})\\
		SFR & Taux de formation stellaire (\textit{star formation rate})\\
		SINGS & Spitzer Infrared Nearby Galaxies Survey\\
		SoFiA & Source Finder Application\\
		THINGS & The HI Nearby Galaxy Survey\\
		THOR & The HI/OH/Recombination Line Survey of the Milky Way\\
		VLA & Very Large Array\\
		WIM & Milieu ionisé tiède (\textit{warm ionized medium})\\
		WNM & Milieu neutre tiède (\textit{warm neutral medium})\\
		\hline\hline
	\end{tabular}
\end{table}
\endgroup
\newpage

\addcontentsline{toc}{chapter}{Nomenclatures}
\section*{{Nomenclatures}}
\begingroup
\renewcommand{\arraystretch}{1.1}
\begin{table}[H]
	\centering
	\begin{tabular}{|c|c|c|}
		\hline
			Symbole & Nomenclature &Valeur / Unité\\
		\hline
			A & Coefficient de normalisation de Kennicutt--Schmidt & - \\
			D & Distance de la galaxie & Mpc \\
			I & Intensité & MJy sr$^{-1}$ \\
			J & Nombre quantique de rotation & - \\
			$\lambda$ & Longueur d'onde & m \\
			$\lambda_A$, $\lambda_\mu$, $\lambda_\sigma$, $\lambda_\sigma''$ & Hyperparamètres de ROHSA & - \\
			N & Indice de la loi de Kennicutt--Schmidt & - \\
			$\nu$ & Fréquence & Hz \\
			$r$ & Rayon & kpc \\
			$r_{25}$ & Rayon optique & kpc \\
			$\rho$ & Densité volumique & g cm$^{-3}$ \\
			$\sigma_\nu$ & Dispersion de vitesse & km s$^{-1}$ \\
			$\Sigma_{  gas}$ & Densité de surface du gaz & M$_\odot$ pc$^{-2}$ \\
			$\Sigma_{  SFR}$ & Densité de surface du taux de formation stellaire & M$_\odot$ kpc$^{-2}$ yr$^{-1}$ \\
			$t_{dep}$ & Temps d'épuisement & yr \\
			$X_{  CO}$ & Facteur de conversion CO--H$_2$ & cm$^{-2}$ (K km s$^{-1}$)$^{-1}$ \\
			$Z$ & Métallicité & $12 + \log(\mathrm{O/H})$ or $0.1\,Z_\odot$ \\
		\hline
	\end{tabular}
\end{table}
\endgroup